Schätzungen: Das Leitungswasser hat etwa $\vartheta_\text{k} = \TkO$, siedendes Wasser $\vartheta_\text{h} = \ThO$, körperwarmes Wasser $\vartheta = \TmO$. Wärmeverluste an Wanne und Luft sowie der kleine Dichteunterschied zwischen heissem und kaltem Wasser sind vernachlässigt. Das heisse Wasser gibt so viel Wärme ab, wie das kalte aufnimmt:
\begin{align}
 c\,\rho V_\text{h}\left(\vartheta_\text{h}-\vartheta\right) &= c\,\rho V_\text{k}\left(\vartheta-\vartheta_\text{k}\right)
\end{align}
Das nötige Volumen heisses Wasser ist
\begin{align}
 V_\text{h} &= \VhF \\
   &= \VkO\cdot\frac{\TmO-\TkO}{\ThO-\TmO} \\
   &= \Vhl \approx \result{\VhlP{0}{1}}
\end{align}
Das Ergebnis hängt stark von der Temperatur des Leitungswassers ab; wegen der Wärmeverluste braucht man in Wirklichkeit eher mehr heisses Wasser.
